\documentclass[11pt]{article}
\usepackage[a4paper,margin=26mm]{geometry}
\usepackage{amsmath,amssymb}
\usepackage[T1]{fontenc}
\usepackage{lmodern,microtype}
\usepackage[colorlinks=true,linkcolor=blue,urlcolor=blue,citecolor=blue]{hyperref}
\hypersetup{pdftitle={Companion Note to "The Unique Games Theorem: A Short Self-Contained Proof": Attribution, Provenance and Checks},pdfauthor={Ariel Elboim}}
\newcommand{\F}{\mathbb F}
\newcommand{\OPT}{\operatorname{OPT}}
\newcommand{\ind}{\mathbf 1}
\newcommand{\eps}{\varepsilon}
\newcommand{\depositdoi}{10.5281/zenodo.23265572}
\title{Companion Note to ``The Unique Games Theorem:\\ A Short Self-Contained Proof'':\\ Attribution, Provenance and Checks}
\author{Ariel Elboim\\\small Independent researcher, ORCID 0009-0001-1341-366X\\\small Prepared with Codex (OpenAI) and Claude (Anthropic)}
\date{9 October 2026}
\begin{document}
\maketitle
\noindent This note accompanies the paper \emph{The Unique Games Theorem:
A Short Self-Contained Proof} \cite{Main} and is deposited with it at
\href{https://doi.org/\depositdoi}{doi:\depositdoi}. The
paper gives a short proof of the theorem of OpenAI's preprint \emph{The
Unique Games Theorem} \cite{UGT}; a longer presentation that follows the
source closely is \cite{Elboim}. Below we list the paper's departures from
the source, credit the contributors to each simplification, describe how
the paper was reviewed, and report a run of the Comparator check on
OpenAI's Lean formalization. Section, theorem, proposition and lemma
numbers refer to the paper.

\section*{Attribution and changes of presentation}

The construction and main proof are from \cite{UGT}. The changes below
simplify the presentation and some proof steps; none is a claim of a new
theorem or an independent verification of the preprint.

\paragraph{Parity equations, tables and weights (Sections 2 and 5).}
The source makes the variables of each equation distinct by replacing
every variable with $100$ clones and listing all distinct-clone choices
with conditional weights; its bound uses sign rounding of clone biases and
conditioning. We instead cancel repeated variables and replace each unary
equation by three equations on a triangle of fresh variables; the bound
$\frac23+\frac13\OPT$ suffices, since only a constant gap is needed, and
turns the source's constants $4/5$, $14/15$ and $3600$ into $6/7$, $20/21$
and $7056$. The source rounds occurrence weights before subdividing; we
keep exact multiplicities, using the equally likely random strings of
H\aa stad's verifier. We define a key as the coarsest factorization of a
table, so pullback compatibility follows from uniqueness; the source
checks that explicit third-bit intercept corrections cancel.

\paragraph{The noise gadget (Section 3).}
The source builds the recursion on translation spaces, starts the
gadget's rank tracking from a generic space, and realizes its output by a
quotient, then copies indexed by injections and a second quotient. In the
paper the recursion is written in direct coordinates
$V_t=W_t\oplus B_{\rm blk}$, so the output is already a subspace and no
quotient is needed; genericity needs only the alignment bound, without
excluding proportional pairs; instead of following one lift (in the
source's sense) through chosen sections, we pull back the whole observer
family; and instead of following the random leaf path with a filtration
and telescoping the expected drops, we prove a detection bound for every
family by induction on the height, whose potential $H_r-\theta N_S/r$ also
admits an arbitrary starting space at a one-time cost below
$2^{r_0}\theta$. The threshold becomes $r_0$, and only output dimensions
divisible by $3$ are realized, which suffices for every later use. In the
kernel identity and the genericity argument, squaring the alignment
equation and formal differentiation over $\F_2(y_{ij})$ replace the
source's square-root monomial-independence and double-annihilator step,
and one annihilating row vector replaces its nonzero augmented minor when
passing to $F$; the closing point-indicator count follows the source. The
source's update space $\{(a,Q(a)+w):a\in A,\ v\cdot w=0\}$ is written in
the paper, via $w\mapsto w+t^2Q(v)$ at $a=tv$, as the kernel of one
$F$-valued equation, which yields its dimension, the update count and the
kernel identity.

\paragraph{Soundness (Section 4).}
The source's matrix-chart argument credits Barak, Kothari and Steurer
\cite{BKS}; the needed calculation is proved in the paper, in the proof
of Lemma~3. The source shows that good row advice is frequent by
recoloring randomly inside $K$ and applying a concentration and counting
bound; we instead bound acceptance directly by $\eta$ plus twice the
measure of the union of good row fibers, which needs neither. The source
selects useful projected data and decodes affine targets $Mz+u$; we keep
constant-output certificates, transfer a score odd along a random nonzero
$h\in\ker A$, and let the second player guess the output and
Fourier-sample one row along $h$, which yields a legal answer at every
sampled frequency; Cauchy--Schwarz then gives unconditional agreement
without revealing the certificate. The noise comparison likewise expands
$\ind_{F_U=u}$ rather than averaging over output characters. The source
takes $k$ along cubes so that $\beta=k^{-2/3}$ is rational; $\beta$ enters
only the soundness experiment, never the output game, so every large $k$
is allowed in the paper. In that experiment, singleton columns of $M_O$
are drawn from $\ker A$ and those of $T_O$ are zero, so every singleton is
clean; the source uses only singletons whose advice vanishes.

\section*{Provenance and review}

Several drafts of the paper received source-based reviews by separate
Claude instances and by Codex; Codex had also helped write an earlier
draft. After reviewing the draft that became the longer presentation
\cite{Elboim}, Codex proposed, at the author's request and in eight
rounds, most of the departures listed under Attribution. Claude proposed
the triangle, which replaced Codex's seven-variable gadget, the keys as
coarsest factorizations, the induction on the height, the update count,
the single annihilator and the update space as one kernel in Section~3,
the agreement bound through one conditional expectation in Section~4, the
two-stage sampling in Section~4.3, and three
restructurings: one soundness proposition (Proposition~2), one lemma from
test acceptance to good row fibers (Lemma~3), and one choice of the gadget
parameters in Section~3. The model gpt-6-astra (OpenAI) proposed the transferred
correlation and a one-row decoder with a random character, and a separate
Codex run through the API (gpt-6.1-sol, OpenAI) proposed a one-row
decoder that guesses the output value instead. Muse Spark 1.3 (Meta)
proposed compressing parts of Sections~2, 3, 4.3 and~5 and reading the
update count through the dot product of the output change with $v$.
Gemini 3.8 Flash (Google) proposed treating the score transfer and the
decoder as one subsection. A separate Claude reviewer checked each round
against the preceding draft, and its fixes are applied. These reviews do
not certify the preprint or constitute formal verification of the paper,
and no human expert has reviewed it.

\section*{Lean check}

OpenAI's repository \cite{UGT} contains a Lean~4 formalization of the
source theorem, with a Comparator challenge that permits only the
standard axioms. We report one run of this check, made on 9 October 2026
at the revision pinned in \cite{UGT}, as an unprivileged user on a
four-core x86-64 cloud machine. It is the run reported in Appendix~A of
\cite{Elboim}. The check script (\texttt{ugc\_lean\_check.py}) and its
complete output logs (\texttt{ugc-lean-ugt-check.log} and the output
directory, as \texttt{ugc\_out.zip}) are deposited with that paper at
\href{https://doi.org/10.5281/zenodo.23264902}{\texttt{doi:10.5281/zenodo.23264902}}.

\begin{center}
\footnotesize
\begin{tabular}{@{}ll@{}}
Repository & \texttt{openai/math}, revision
 \texttt{fd4aeeb2ee4fc729c18d98444fed42fd0529eeeb}\\
Tools & Lean 4.34.1 (pinned toolchain); Comparator commit \texttt{d03acab}
 (tag \texttt{v4.34.0})\\
Command & \texttt{lake env comparator
 ComparatorChallenges/UniqueGamesTheorem.json}\\
Declaration & \texttt{OAI.UniqueGamesTheorem.theorem11}\\
Axioms & \texttt{propext}, \texttt{Classical.choice},
 \texttt{Quot.sound}\\
Verdict & ``Lean default kernel accepts the solution'', ``Your solution
 is okay!''
\end{tabular}
\end{center}

\paragraph{Procedure and result.} Following the repository's
instructions, our script ran \texttt{lake update} (manifest and toolchain
unchanged), \texttt{lake exe cache get} and the command above in the
directory \texttt{lean}, then a separate \texttt{lake build} of the
solution and challenge modules and \texttt{\#print axioms}. The challenge
module states the declaration's type: \texttt{BinaryGapReduction}~$\eps$~$\delta$
is nonempty for all real $\eps,\delta\in(0,1/2)$. Comparator exited with
code~$0$, the build succeeded, and the solution's axioms include no
\texttt{sorryAx}; only the challenge module uses \texttt{sorry}, by design,
to state the type. The run took 66 minutes.

\paragraph{Caveats.} The sandbox was not enforced: \texttt{landrun}
started, but in our test a write outside its allowed directory succeeded,
and the \texttt{systemd-run} guard recommended in Comparator's README was
unavailable. The check therefore lacked the recommended isolation. We did
not audit whether the formal statement matches Theorem~1 of the paper.
Above all, the run concerns OpenAI's formal development of the source
theorem. It does not check the paper's simplifications listed under
Attribution, or anything else written in the paper, and the paper is not
derived from the formalization.

{\small
\begin{thebibliography}{9}
\bibitem{Main}
Ariel Elboim. \emph{The Unique Games Theorem: A Short Self-Contained
Proof}. 2026. Deposited with this note at
\href{https://doi.org/\depositdoi}{doi:\depositdoi}.

\bibitem{UGT}
OpenAI. \emph{The Unique Games Theorem}. Preprint, 23 September 2026,
58 pages. Pinned source revision:
\href{https://github.com/openai/math/tree/fd4aeeb2ee4fc729c18d98444fed42fd0529eeeb/preprints/The-Unique-Games-Theorem-September-23-2026}{GitHub source and paper}.

\bibitem{Elboim}
Ariel Elboim. \emph{A Compressed Presentation of OpenAI's Proof of the
Unique Games Theorem}. Preprint, 2026,
\href{https://doi.org/10.5281/zenodo.23264902}{doi:10.5281/zenodo.23264902}.

\bibitem{BKS}
Boaz Barak, Pravesh K. Kothari, and David Steurer.
\emph{Small-Set Expansion in Shortcode Graph and the 2-to-2 Conjecture}.
Proceedings of ITCS, LIPIcs 124, 9:1--9:12, 2019.
\end{thebibliography}
}
\end{document}
